% ===================================================================
% Exercises -- Introduction to Communication Networks (HUJI 67594)
% One block per week: questions first, then full solutions on a new page.
% Sources: [KR x.Py] = Kurose & Ross 8th ed., chapter x, problem y (reworded);
% [Rn]/[Tn] = recitation/tutorial n (variant with new numbers); [own] = new.
% Update every week: add \weekbar, a questions block and a solutions block.
% Build: pdflatex -interaction=nonstopmode exercises.tex
% ===================================================================
% ===================================================================
% Preamble for the Introduction to Communication Networks weekly notes (copy of the shared
% lecture-notes preamble, kept here so this folder builds on its own).
% ===================================================================
\documentclass[11pt,a4paper]{article}

% ---------- encoding & fonts ----------
\usepackage[T1]{fontenc}
\usepackage[utf8]{inputenc}
\usepackage{lmodern}
\renewcommand{\familydefault}{\rmdefault}
\usepackage{microtype}

% ---------- layout ----------
\usepackage[a4paper,margin=2.3cm,top=2.4cm,bottom=2.6cm,headsep=0.9cm]{geometry}
\setlength{\parindent}{0pt}
\setlength{\parskip}{0.55em}
\usepackage{needspace}
\usepackage{multicol}
\usepackage[shortlabels]{enumitem}
\usepackage{array,booktabs,colortbl,tabularx}
\newcolumntype{L}[1]{>{\raggedright\arraybackslash}p{#1}}

% ---------- maths ----------
\usepackage{amsmath,amssymb,mathtools}

% ---------- colours ----------
\usepackage[dvipsnames]{xcolor}
\definecolor{paper}{HTML}{EFE9DA}
\definecolor{surface}{HTML}{FDFBF4}
\definecolor{surface2}{HTML}{F1EBDB}
\definecolor{ink}{HTML}{221C35}
\definecolor{muted}{HTML}{6E6883}
\definecolor{border}{HTML}{4B3C8E}
\definecolor{shadowcol}{HTML}{2E2552}
\definecolor{pgreen}{HTML}{1D6A3B}
\definecolor{ppink}{HTML}{D0307A}
\definecolor{pgold}{HTML}{8C5C0B}
\definecolor{pcyan}{HTML}{2B5FA8}
\pagecolor{paper}
\color{ink}

% ---------- graphics ----------
\usepackage{tikz}
\usetikzlibrary{positioning,arrows.meta,calc,shapes.geometric,shapes.misc,decorations.pathreplacing,patterns,fit,backgrounds}
\usepackage{pgfplots}
\pgfplotsset{compat=1.18}
\tikzset{
  lbl/.style={font=\footnotesize,text=ink,inner sep=1pt},
  blk/.style={draw=ink,line width=1pt,fill=surface,font=\sffamily\scriptsize\bfseries,align=center,minimum height=0.85cm,inner sep=4pt},
  arr/.style={-{Stealth},line width=1.1pt,ink},
  dot/.style={circle,fill=ink,inner sep=1.4pt},
  plbl/.style={font=\ttfamily\scriptsize,text=muted,inner sep=1pt},
  >={Stealth[length=2.6mm]}
}

% ---------- inline semantic markup ----------
\newcommand{\term}[1]{\textcolor{pgreen}{#1}}   % newly defined term
\newcommand{\key}[1]{\textcolor{pcyan}{#1}}     % key word / unit
\newcommand{\hl}[1]{\textcolor{ppink}{#1}}      % the thing to remember
\newcommand{\thd}[1]{{\color{surface}\sffamily\bfseries\footnotesize #1}} % table header cell
\newcommand{\src}[1]{{\ttfamily\scriptsize\textcolor{muted}{[#1]}}}

% ---------- boxes ----------
\usepackage[most]{tcolorbox}
\tcbset{
  notebase/.style={
    enhanced jigsaw, breakable, sharp corners,
    colback=surface, colframe=border, boxrule=1.1pt,
    left=6pt,right=6pt,top=5pt,bottom=5pt,
    before skip=1.1em, after skip=1.1em,
    shadow={1.4mm}{-1.4mm}{0mm}{fill=shadowcol},
    fonttitle=\ttfamily\scriptsize,
    attach boxed title to top left={xshift=3mm,yshift*=-2.4mm},
    boxed title style={sharp corners,boxrule=0.9pt,colframe=surface,
      left=4pt,right=4pt,top=2.5pt,bottom=2.5pt,
      shadow={0.8mm}{-0.8mm}{0mm}{fill=shadowcol}},
    parbox=false,
    before upper={\setlength{\parskip}{0.45em}},
  }
}
\newcommand{\boxtitle}[2]{\if\relax\detokenize{#2}\relax #1\else #1 \textperiodcentered\ \MakeUppercase{#2}\fi}
\newcommand{\boxtitleonly}[2]{\if\relax\detokenize{#2}\relax #1\else \MakeUppercase{#2}\fi}

\newtcolorbox{defn}[1][]{notebase,colbacktitle=pgreen,title={\boxtitle{DEF}{#1}}}
\newtcolorbox{formula}[1][]{notebase,colbacktitle=border,title={\boxtitleonly{FORMULA}{#1}}}
\newtcolorbox{intuition}[1][]{notebase,colbacktitle=ppink,title={\boxtitleonly{KEY IDEA}{#1}}}
\newtcolorbox{example}[1][]{notebase,colbacktitle=pcyan,title={\boxtitle{EXAMPLE}{#1}}}
\newtcolorbox{pitfall}[1][]{notebase,colbacktitle=pgold,title={\boxtitle{PITFALL}{#1}}}
\newtcolorbox{plainbox}{notebase}

% lists inside notes
\setlist[itemize]{label={\textcolor{pgreen}{\scriptsize$\blacksquare$}},leftmargin=1.3em,itemsep=0.25em,topsep=0.3em}
\setlist[enumerate]{label={\textcolor{pgreen}{\ttfamily\arabic*.}},leftmargin=1.6em,itemsep=0.3em,topsep=0.3em}

% ---------- headings ----------
\usepackage{titlesec}
\newcommand{\twodigit}[1]{\ifnum#1<10 0\fi\number#1}
\titleformat{\section}[display]
  {\sffamily\bfseries\fontsize{20}{24}\selectfont\color{pgold}}
  {\ttfamily\small\mdseries\color{pgreen}SECTION \twodigit{\value{section}}}
  {8pt}{}[\vspace{4pt}{\color{border}\titlerule[1.4pt]}]
\titlespacing*{\section}{0pt}{1.6em}{1.0em}
\titleformat{\subsection}
  {\ttfamily\large\color{pcyan}}{>\,\thesubsection}{0.6em}{}
\titlespacing*{\subsection}{0pt}{1.3em}{0.5em}
\titleformat{\subsubsection}
  {\sffamily\bfseries\color{border}}{\thesubsubsection}{0.6em}{}
\titlespacing*{\subsubsection}{0pt}{1.0em}{0.3em}

% ---------- running header / footer ----------
\usepackage{fancyhdr}
\newcommand{\headleft}{}
\newcommand{\headright}{}
\pagestyle{fancy}
\fancyhf{}
\renewcommand{\headrulewidth}{0pt}
\fancyhead[L]{\ttfamily\scriptsize\textcolor{muted}{\headleft}}
\fancyhead[R]{\ttfamily\scriptsize\textcolor{muted}{\headright}}
\fancyfoot[C]{\tikz{\node[draw=border,line width=1pt,fill=surface,inner sep=2.5pt,font=\sffamily\footnotesize] (p) {\thepage};
  \draw[shadowcol,line width=1.6pt] ([yshift=-2pt]p.south west)--([yshift=-2pt]p.south east);}}
\fancypagestyle{plain}{\fancyhf{}\renewcommand{\headrulewidth}{0pt}\fancyfoot[C]{\tikz{\node[draw=border,line width=1pt,fill=surface,inner sep=2.5pt,font=\sffamily\footnotesize] (p) {\thepage};
  \draw[shadowcol,line width=1.6pt] ([yshift=-2pt]p.south west)--([yshift=-2pt]p.south east);}}}

% ---------- links ----------
\usepackage[hidelinks]{hyperref}

\newgeometry{a4paper,margin=1.8cm,top=2cm,bottom=2cm,headsep=0.6cm}
\renewcommand{\headleft}{67594 \textperiodcentered\ EXERCISES}
\renewcommand{\headright}{INTRODUCTION TO COMMUNICATION NETWORKS}
\setlength{\parskip}{0.3em}

\newcommand{\E}{\mathbb{E}}
\newcommand{\PROP}{\text{PROP}}
\newcommand{\TRANSP}{\text{TRANSP}}
\newcommand{\cmpl}[1]{\overline{#1}}

\titleformat{\section}{\sffamily\bfseries\large\color{pgold}}{}{0pt}{}[\vspace{-2pt}{\color{border}\titlerule[0.8pt]}]
\titlespacing*{\section}{0pt}{0.9em}{0.5em}
\pretocmd{\section}{\needspace{6\baselineskip}}{}{}

\newcommand{\weekbar}[3][\clearpage]{%
  #1\noindent\colorbox{border}{\parbox{\dimexpr\linewidth-2\fboxsep}{%
  \color{surface}\sffamily\bfseries\large WEEK #2\hfill{\normalsize\mdseries #3}}}\par\medskip}
\newcommand{\solbar}[1]{%
  \clearpage\noindent\colorbox{pgreen}{\parbox{\dimexpr\linewidth-2\fboxsep}{%
  \color{surface}\sffamily\bfseries\large WEEK #1 \textperiodcentered\ SOLUTIONS}}\par\medskip}

% question / solution headers: \q{number}{difficulty 1-3}{source}
\newcommand{\diff}[1]{\ifcase#1\or$\bullet\circ\circ$\or$\bullet\bullet\circ$\or$\bullet\bullet\bullet$\fi}
\newcommand{\q}[3]{\par\needspace{3\baselineskip}\smallskip\noindent{\sffamily\bfseries\color{border}Q#1}\quad{\color{pgold}\diff{#2}}\quad\src{#3}\par\nopagebreak}
\newcommand{\s}[1]{\par\needspace{3\baselineskip}\smallskip\noindent{\sffamily\bfseries\color{pgreen}Q#1}\par\nopagebreak}

\newcommand{\idea}[1]{\par\noindent{\color{pcyan}\sffamily\bfseries Idea.}\ #1\par}
\newcommand{\ans}[1]{\par\noindent\colorbox{surface2}{\parbox{\dimexpr\linewidth-2\fboxsep}{{\color{pgreen}\sffamily\bfseries Answer.}\ #1}}\par}
\newcommand{\watch}[1]{\par\noindent{\color{pgold}\sffamily\bfseries Watch out:}\ #1\par}

\begin{document}

{\sffamily\bfseries\fontsize{24}{28}\selectfont\color{pgold}Exercises}\hfill{\ttfamily\small\color{pgreen}HUJI 67594 \textperiodcentered\ updated weekly}\par
\vspace{2pt}{\color{border}\rule{\linewidth}{1.4pt}}\par
\begin{plainbox}
\textbf{How to use.} Each week has its questions first and its solutions on a separate page, so you can work without seeing the answers. Every solution gives the \emph{idea} (which tool and why), numbered steps with all numbers, a boxed \emph{answer}, and a \emph{watch out} line where people often slip.
Difficulty: \diff1 direct use of a formula, \diff2 several steps, \diff3 exam level. Sources:
\src{KR 6.P11} = Kurose \& Ross, 8th ed., chapter 6, problem 11 (reworded); \src{R2}, \src{T1} = a variant of a recitation or tutorial question with new numbers; \src{own} = written for this sheet.
The formulas are in \texttt{formula\_sheet.pdf}; the explanations are in \texttt{weekNN.pdf}.
\end{plainbox}

\noindent\textbf{Contents.}\quad Week 01: delay, switching, probability (14 questions). \quad Week 02: layering, ALOHA, CSMA/CD, WiFi (18 questions).

% =================================================================
\weekbar[\medskip]{01}{Lecture 1: Introduction \textperiodcentered\ Tutorial 1: Probability}
% =================================================================

\section{Delay and switching}

\q{1.1}{1}{KR 1.P6}
Hosts A and B are $m$ metres apart, joined by one link of rate $R$ bit/s with propagation speed $s$. A sends one packet of $L$ bits starting at $t=0$.
(a) Write $d_{\text{prop}}$, $d_{\text{trans}}$ and the end-to-end delay (no processing or queuing).
(b) At $t=d_{\text{trans}}$, where is the last bit?
(c) If $d_{\text{prop}}>d_{\text{trans}}$, where is the first bit at $t=d_{\text{trans}}$?
(d) With $s=2.5\times10^8$\,m/s, $L=1500$ bytes and $R=10$\,Mbps, for which $m$ is $d_{\text{prop}}=d_{\text{trans}}$?

\q{1.2}{2}{KR 1.P10}
A 1500-byte packet crosses three links joined by two switches. Every link has $R=2.5$\,Mbps and $s=2.5\times10^8$\,m/s; the links are 5000, 4000 and 1000\,km long.
Each switch adds a processing delay of 3\,ms. There is no queuing. Find the end-to-end delay.

\q{1.3}{1}{KR 1.P12}
A packet arrives at a router. On its outgoing link one packet is half transmitted and four more are waiting. All packets are 1500 bytes and the link is 2.5\,Mbps.
(a) What is the queuing delay of the new packet? (b) Give the general formula when $x$ bits of the current packet are already sent and $n$ packets wait.

\q{1.4}{2}{KR 1.P31}
A message of $10^6$ bits goes from a source through two store-and-forward switches to a destination (three links, 5\,Mbps each). Ignore propagation, queuing and processing.
(a) How long does it take without segmentation?
(b) Now it is cut into 100 packets of 10\,000 bits. When does the first packet reach the first switch? When is the second packet fully received there?
(c) How long does the whole message take now? Why is it faster?

\q{1.5}{1}{own}
A 1\,000\,000-bit file is sent over a circuit-switched path. Every link is a 2\,Mbps TDM link with 20 slots, and setting up the circuit takes 300\,ms.
How long until the whole file has arrived (ignore propagation)? Does it depend on the number of links?

\q{1.6}{1}{own}
An application has peak rate 1\,Mbps and average rate 50\,kbps. (a) What is its utilization under circuit switching? (b) How many such applications fit on a 10\,Mbps link
with circuit switching? (c) How many would fit if each needed only its average rate?

\q{1.7}{1}{own}
Packets of 1000 bits arrive at a 1\,Mbps link at $a=900$ packets/s on average. (a) What is the traffic intensity? (b) Up to which arrival rate is the queue stable?

\q{1.8}{2}{KR 1.P8}
Users share a 10\,Mbps link. Each needs 200\,kbps when active and is active 10\% of the time, independently.
(a) How many users does circuit switching support?
(b) With packet switching and 120 users, what is the probability that exactly $n$ users are active?
(c) What is the probability that 51 or more are active? What does that mean for packet switching?

\q{1.9}{2}{KR 1.P9}
Same as Q1.8 but with a 1\,Gbps link and users of 100\,kbps, each active with probability $p=0.1$.
(a) What is $N$, the number of users under circuit switching? (b) With packet switching and $M$ users, write the probability that more than $N$ are active.

\section{Probability}

\q{1.10}{2}{T1 / lecture Example II}
Good packets survive $t$ seconds with probability $e^{-t}$, bad ones with probability $e^{-1000t}$. A packet is good with probability $0.2$.
(a) What is the probability that a packet survives $2$\,ms? (b) It survived $2$\,ms: what is the probability that it is good?

\q{1.11}{1}{KR 1.P22}
A path has 5 links, each losing a packet independently with probability $0.05$. The server resends a lost packet end to end.
(a) Probability that one transmission arrives? (b) Expected number of \emph{re}transmissions?

\q{1.12}{3}{T1 / lecture Example III}
A chain of $n=4$ links, each failing independently with probability $p=0.2$.
(a) With end-to-end retransmission, what is the expected number of attempts by S, and the expected number of \emph{link} transmissions?
(b) With hop-by-hop retransmission, what is the expected number of link transmissions? Which scheme is cheaper?

\q{1.13}{2}{own}
Packets arrive as a Poisson process with rate $\lambda=3$ per second.
(a) Probability of no packet in 1\,s? (b) Probability of exactly 2 packets in 0.5\,s?
(c) You have already waited 2\,s without a packet. What is the expected remaining wait? Why?

\q{1.14}{1}{T1}
In each interval of length $T$ the number of packets sent is $\text{Pois}(gT)$. (a) What is the probability of exactly one packet? (b) For which $gT$ is it largest, and what is the largest value?

% =================================================================
\solbar{01}
% =================================================================

\s{1.1}
\idea{Two different delays. Transmission $=$ the time to push all the bits out of A (depends on $L$ and $R$). Propagation $=$ the time one bit needs to travel the wire (depends on $m$ and $s$).}
\begin{enumerate}
  \item[(a)] $d_{\text{prop}}=m/s$ and $d_{\text{trans}}=L/R$. The last bit leaves A at $t=L/R$ and needs $m/s$ more to arrive, so the end-to-end delay is $L/R+m/s$.
  \item[(b)] At $t=d_{\text{trans}}$, A has just pushed out its last bit: the last bit is \emph{at A}, just leaving.
  \item[(c)] The first bit left at $t=0$ and has travelled for $d_{\text{trans}}$ seconds, i.e.\ $s\cdot d_{\text{trans}}$ metres. Since $d_{\text{prop}}>d_{\text{trans}}$, this is less than $m$: the first bit is still on the wire and has not reached B.
  \item[(d)] Set $m/s=L/R$, so $m=sL/R$. $L=1500\times8=12\,000$ bits, $L/R=12\,000/10^7=1.2$\,ms, $m=2.5\times10^8\times1.2\times10^{-3}=3\times10^5$\,m.
\end{enumerate}
\ans{(a) $L/R+m/s$; (b) at A, leaving; (c) $s\,d_{\text{trans}}$ from A, still on the wire; (d) $m=300$\,km.}
\watch{bytes $\to$ bits ($\times8$) before dividing by $R$.}

\s{1.2}
\idea{Store-and-forward: each switch waits for the \emph{whole} packet before sending it on, so the transmission time is paid on every link. Add the four kinds of delay.}
\begin{enumerate}
  \item Transmission: $L=1500\times8=12\,000$ bits; per link $12\,000/(2.5\times10^6)=4.8$\,ms; three links: $14.4$\,ms.
  \item Propagation: $5\times10^6/2.5\times10^8=20$\,ms, $4\times10^6/2.5\times10^8=16$\,ms, $10^6/2.5\times10^8=4$\,ms; total $40$\,ms.
  \item Processing: only at the 2 switches: $2\times3=6$\,ms.
  \item Queuing: $0$ (given).
\end{enumerate}
\ans{$14.4+40+6=60.4$\,ms.}
\watch{km $\to$ m; processing happens at 2 switches, not 3.}

\s{1.3}
\idea{The new packet waits until every bit ahead of it has been transmitted.}
\begin{enumerate}
  \item One packet takes $L/R=12\,000/(2.5\times10^6)=4.8$\,ms.
  \item Ahead of it: the remaining half of the current packet ($2.4$\,ms) and 4 full packets ($4\times4.8=19.2$\,ms).
  \item General case: the current packet still has $L-x$ bits to send, and $n$ packets hold $nL$ bits.
\end{enumerate}
\ans{(a) $2.4+19.2=21.6$\,ms; (b) $\dfrac{nL+(L-x)}{R}$.}

\s{1.4}
\idea{Without segmentation every link must receive the whole message before the next link can start. With packets, the links work at the same time, like a pipeline.}
\begin{enumerate}
  \item[(a)] The message takes $10^6/(5\times10^6)=0.2$\,s on one link. It is stored at each switch, so the 3 links are used one after the other: $3\times0.2=0.6$\,s.
  \item[(b)] One packet takes $10^4/(5\times10^6)=2$\,ms per link. Packet 1 reaches switch 1 at $t=2$\,ms. From $2$ to $4$\,ms the source sends packet 2 (while switch 1 forwards packet 1), so packet 2 is fully at switch 1 at $t=4$\,ms.
  \item[(c)] Packet 1 needs $3\times2=6$\,ms to reach the destination. After that, one packet arrives every $2$\,ms, so the other 99 arrive $99\times2=198$\,ms later. Total $6+198=204$\,ms. This is the formula $(N+K-1)L/R=(3+100-1)\times2$\,ms.
\end{enumerate}
\ans{(a) $0.6$\,s; (b) $2$\,ms and $4$\,ms; (c) $204$\,ms, about 3 times faster, because all three links are busy at the same time.}

\s{1.5}
\idea{TDM gives the circuit a fixed share of the link rate. The setup is paid once, and then the bits stream through.}
\begin{enumerate}
  \item The circuit's rate is $2\,\text{Mbps}/20=100$\,kbps.
  \item Transmission: $10^6/10^5=10$\,s.
  \item Add the setup: $10+0.3=10.3$\,s.
\end{enumerate}
\ans{$10.3$\,s. It does not depend on the number of links: the switches do not store the file, they pass the bits straight through the reserved path.}

\s{1.6}
\idea{A circuit must reserve the \emph{peak} rate. Utilization $=$ what is used on average $\div$ what is reserved.}
\begin{enumerate}
  \item[(a)] $A/P=50\,\text{kbps}/1000\,\text{kbps}=5\%$.
  \item[(b)] Each reserves $1$\,Mbps: $10/1=10$ applications.
  \item[(c)] Each needs only $50$\,kbps: $10\,000/50=200$ applications.
\end{enumerate}
\ans{5\%; 10; 200. Packet switching (statistical multiplexing) gets close to (c), because the applications rarely peak at the same time.}

\s{1.7}
\idea{Traffic intensity $I=La/R$ $=$ bits arriving per second $\div$ bits the link can send per second.}
\begin{enumerate}
  \item[(a)] Arriving: $1000\times900=9\times10^5$ bit/s. $I=9\times10^5/10^6=0.9$.
  \item[(b)] Stable needs $La<R$: $a<10^6/1000=1000$ packets/s.
\end{enumerate}
\ans{(a) $0.9$; (b) $a<1000$ packets/s.}
\watch{$I<1$ only means the queue does not grow forever. At $I=0.9$ the queuing delay is already large.}

\s{1.8}
\idea{Circuit switching: divide the link rate. Packet switching: count the active users, which is a binomial random variable.}
\begin{enumerate}
  \item[(a)] $10\,\text{Mbps}/200\,\text{kbps}=50$ users.
  \item[(b)] Each of the 120 users is active independently with probability $0.1$, so $X\sim\text{Bin}(120,0.1)$ and $P(X=n)=\binom{120}{n}0.1^n0.9^{120-n}$.
  \item[(c)] The link holds 50 active users, so overload means $X\ge51$: $P(X\ge51)=\sum_{n=51}^{120}\binom{120}{n}0.1^n0.9^{120-n}$. Sense check: the mean is $np=12$ and the standard deviation $\sqrt{np(1-p)}=\sqrt{10.8}\approx3.3$, so 51 is $(51-12)/3.3\approx12$ standard deviations above the mean. Numerically the sum is about $2\times10^{-20}$.
\end{enumerate}
\ans{(a) 50; (b) $\binom{120}{n}0.1^n0.9^{120-n}$; (c) $\approx2\times10^{-20}$: packet switching serves 120 users where circuits serve 50, and is practically never overloaded.}

\s{1.9}
\begin{enumerate}
  \item[(a)] $N=10^9/10^5=10\,000$ users.
  \item[(b)] Active users $X\sim\text{Bin}(M,p)$. ``More than $N$'' means $X\ge N+1$.
\end{enumerate}
\ans{(a) $10\,000$; (b) $P(X>N)=\sum_{k=N+1}^{M}\binom Mk p^k(1-p)^{M-k}=1-\sum_{k=0}^{N}\binom Mk p^k(1-p)^{M-k}$.}

\s{1.10}
\idea{(a) We do not know the packet's type, so split on it (total probability). (b) We know the result and ask about the cause, so reverse the condition (Bayes).}
Let $A$ = ``survives $t=2$\,ms'' and $B$ = ``good''. Given: $P(B)=0.2$, $P(A\mid B)=e^{-0.002}=0.998$, $P(A\mid\cmpl B)=e^{-1000\cdot0.002}=e^{-2}=0.135$.
\begin{enumerate}
  \item[(a)] $P(A)=P(A\mid B)P(B)+P(A\mid\cmpl B)P(\cmpl B)=0.998\cdot0.2+0.135\cdot0.8=0.1996+0.1083=0.308$.
  \item[(b)] $P(B\mid A)=\dfrac{P(A\mid B)P(B)}{P(A)}=\dfrac{0.1996}{0.308}=0.648$.
\end{enumerate}
\ans{(a) $0.308$; (b) $0.648$. Surviving raised the chance of ``good'' from $0.2$ to $0.65$, because bad packets rarely survive even 2\,ms.}

\s{1.11}
\begin{enumerate}
  \item[(a)] The packet must survive all 5 links; they are independent, so multiply: $0.95^5=0.774$.
  \item[(b)] Each attempt succeeds independently with probability $0.774$, so the number of attempts is $\text{Geo}(0.774)$ with mean $1/0.774=1.292$. The first attempt is not a retransmission: $1.292-1=0.292$.
\end{enumerate}
\ans{(a) $0.774$; (b) $0.29$ retransmissions on average.}
\watch{``retransmissions'' $=$ attempts $-1$.}

\s{1.12}
\idea{End to end: a loss anywhere repeats the whole path. Hop by hop: a loss repeats only that link. Count link transmissions in both cases.}
\begin{enumerate}
  \item[(a)] One attempt succeeds with probability $0.8^4=0.4096$, so the attempts are $\text{Geo}(0.4096)$: mean $1/0.4096=2.44$.
  An attempt always uses link 1; it uses link 2 only if link 1 worked ($0.8$), link 3 only if links 1--2 worked ($0.64$), link 4 with probability $0.512$.
  By linearity of expectation, one attempt uses $1+0.8+0.64+0.512=2.952$ links on average. Total: $2.44\times2.952\approx7.2$ link transmissions.
  \item[(b)] Each hop is retried until it works: $\text{Geo}(0.8)$, mean $1/0.8=1.25$ sends. Four hops: $4\times1.25=5$.
\end{enumerate}
\ans{(a) $2.44$ attempts, $\approx7.2$ link transmissions; (b) $5$. Hop by hop is cheaper.}

\s{1.13}
\idea{In a Poisson process with rate $\lambda$, the count in a window of length $t$ is $\text{Pois}(\lambda t)$, and the gaps are $\text{Exp}(\lambda)$, which is memoryless.}
\begin{enumerate}
  \item[(a)] 1\,s: $\text{Pois}(3)$. $P(0)=e^{-3}=0.050$.
  \item[(b)] 0.5\,s: $\text{Pois}(1.5)$. $P(2)=e^{-1.5}\,\dfrac{1.5^2}{2!}=0.2231\times1.125=0.251$.
  \item[(c)] Memorylessness: the remaining wait is again $\text{Exp}(3)$, with mean $1/3$\,s, whatever the time already waited.
\end{enumerate}
\ans{(a) $0.050$; (b) $0.251$; (c) $1/3$\,s.}
\watch{Scale the mean with the window: $\lambda t$, not $\lambda$.}

\s{1.14}
\begin{enumerate}
  \item[(a)] $P(X=1)=\dfrac{(gT)^1}{1!}e^{-gT}=gT\,e^{-gT}$.
  \item[(b)] Write $x=gT$ and $f(x)=xe^{-x}$. Then $f'(x)=e^{-x}-xe^{-x}=(1-x)e^{-x}$, which is zero at $x=1$; $f$ rises before and falls after, so it is a maximum. $f(1)=e^{-1}$.
\end{enumerate}
\ans{(a) $gT\,e^{-gT}$; (b) largest at $gT=1$, value $1/e\approx0.368$. This is slotted ALOHA's best goodput (week 2).}

% =================================================================
\weekbar{02}{Lecture 2: Layering \& Broadcast \textperiodcentered\ Recitation 2: ALOHA}
% =================================================================

\section{Layering}

\q{2.1}{1}{own}
An application sends 200 bytes over TCP/IP/Ethernet. (a) What fraction of the headers-plus-data is useful data? (b) What if you also count Ethernet's 4-byte CRC and 8-byte preamble?

\q{2.2}{1}{R2}
Give the OSI layer of: HTTP, TLS encryption, TCP acknowledgements, IP routing, CSMA/CD, the optical fibre. Which of them are not a separate layer in the Internet's 5-layer stack?

\section{Channel partitioning and ALOHA}

\q{2.3}{1}{own}
A 40\,Mbps channel is shared by TDMA among 8 stations; only 3 have data. (a) Rate of each busy station? (b) How much capacity is wasted? (c) Which wish-list property fails?

\q{2.4}{2}{KR 6.P8, P9}
(a) For slotted ALOHA with $N$ nodes, find the $p$ that maximizes $Np(1-p)^{N-1}$ and the limit of the maximum as $N\to\infty$.
(b) Do the same for pure ALOHA, $Np(1-p)^{2(N-1)}$.

\q{2.5}{2}{KR 6.P11}
Four nodes A, B, C, D use slotted ALOHA; each always has frames and sends in each slot with probability $p$. Slots are numbered $1,2,\dots$
(a) Probability that A succeeds for the first time in slot 4?
(b) Probability that some node succeeds in slot 5?
(c) Probability that the first success (by anyone) is in slot 4?
(d) Efficiency of the system? Give all answers as formulas and as numbers for $p=0.25$.

\q{2.6}{3}{KR 6.P10}
Two nodes A and B use slotted ALOHA with probabilities $p_A$ and $p_B$.
(a) Give A's average throughput (fraction of slots A wins) and the total efficiency.
(b) If $p_A=2p_B$, is A's throughput twice B's? If not, which $p_A$ makes it exactly twice?
(c) With $N$ nodes, A using $2p$ and all others $p$: the throughput of A and of any other node?

\q{2.7}{2}{R2}
A slotted ALOHA channel has $B=2$\,Mbps, frames of $L=2000$ bits, and $n=8$ nodes, each sending with $p=0.1$.
(a) Slot length? (b) Goodput? (c) Throughput (fraction of slots with any transmission)? (d) Useful bit rate? (e) Average time between successful frames?

\q{2.8}{2}{R2 / own}
A channel carries on average $G=0.4$ attempts per slot (Poisson model).
(a) Slotted ALOHA: fraction of slots that are successful, empty, and collisions? (b) Pure ALOHA goodput? (c) Should the load be increased or decreased to maximize each?

\q{2.9}{3}{R2 Q1 variant}
Hosts C and D are both attached to two networks. The upper one runs slotted ALOHA and has 3 other nodes; the lower one runs pure ALOHA and has 2 other nodes. D never sends.
Every other node sends in each of its slots with probability $q$. When C has a message for D it picks the upper network with probability $p$ and then sends with probability $q$.
(a) C has one attempt: probability that the message arrives? Numbers for $p=0.6$, $q=0.1$.
(b) C now tries at the start of every slot (choosing a network each time). Goodput of the upper network? Same numbers.

\q{2.10}{3}{R2 Q2 variant}
One channel carries $n$ slotted-ALOHA nodes (slots of length $T$, aligned to a common grid) and $m$ pure-ALOHA nodes (frames of length $T$, own clocks, not aligned with the grid).
Each node sends with probability $p$ at the start of each of its slots.
(a) Probability that a given slotted node succeeds in a slot? (b) That a given pure node succeeds? (c) Numbers for $n=m=5$, $p=0.1$. Who does better, and why?

\section{CSMA/CD, Ethernet, WiFi}

\q{2.11}{1}{own}
A CSMA/CD LAN runs at 1\,Gbps over at most 200\,m of cable ($s=2\times10^8$\,m/s). What is the minimum frame size? Compare with Ethernet's 64 bytes.

\q{2.12}{1}{KR 6.R6}
After the 5th collision, what is the probability that a node picks $K=4$? How long is that delay on 10\,Mbps Ethernet?

\q{2.13}{1}{KR 6.P17}
For $K=100$, how long does an adapter wait on a 100\,Mbps channel? On a 1\,Gbps channel?

\q{2.14}{2}{KR 6.P18}
A and B are on a 10\,Mbps channel, 325 bit times apart. A starts sending a minimum-size frame ($512+64=576$ bit times including the preamble) at $t=0$; B starts before A's signal reaches it.
Can A finish its frame before it hears B? What does that mean?

\q{2.15}{2}{own}
Two NICs have just collided for the 2nd time. (a) From which set is $K$ chosen? (b) Probability that they collide again? (c) Probability they collide on each of the next three attempts (3rd, 4th and 5th)?

\q{2.16}{2}{lecture}
On a CSMA/CD LAN, $\PROP=2\,\mu$s and $\TRANSP=50\,\mu$s.
(a) Efficiency in the simple model ($A=1.5$) and the precise model? (b) The simple model with the large-$N$ value $\alpha=1/e$ instead of $0.4$?

\q{2.17}{3}{KR 6.P20}
A CSMA/CD-like protocol: time is slotted with short slots; a frame lasts $k$ slots. When the channel is free, each of $N$ nodes transmits in a slot with probability $p$.
If exactly one does, it holds the channel for its whole frame ($k$ slots); otherwise the slot is wasted.
(a) Efficiency for fixed $N$, $p$? (b) The best $p$? (c) The efficiency as $N\to\infty$? (d) What happens as $k\to\infty$?

\q{2.18}{1}{own}
In 802.11, RTS and CTS take $25\,\mu$s each and a data frame takes $500\,\mu$s.
(a) What is the reservation overhead per frame? (b) How much time does a collision waste with and without RTS/CTS? (c) When is RTS/CTS not worth it?

% =================================================================
\solbar{02}
% =================================================================

\s{2.1}
\idea{Every packet carries $20$ (TCP) $+\,20$ (IP) $+\,14$ (Ethernet) $=54$ bytes of headers.}
\begin{enumerate}
  \item[(a)] Total $200+54=254$ bytes; useful $200/254$.
  \item[(b)] Add the 4-byte CRC and the 8-byte preamble: $254+12=266$ bytes; useful $200/266$.
\end{enumerate}
\ans{(a) $78.7\%$; (b) $75.2\%$.}

\s{2.2}
\idea{Ask what each thing does: meaning for the user (7), data format or encryption (6), process-to-process delivery (4), finding a path between networks (3), sharing one link (2), the medium (1).}
\begin{itemize}
  \item HTTP: 7. TLS: 6. TCP ACKs: 4. IP routing: 3. CSMA/CD: 2 (the MAC part of layer 2). Fibre: 1.
\end{itemize}
\ans{as above. The Internet's 5-layer stack has no separate session (5) or presentation (6) layer: the application does that work itself, so TLS counts as ``application''.}

\s{2.3}
\idea{TDMA gives every station a fixed $1/8$ of the time, busy or not.}
\begin{enumerate}
  \item[(a)] $40/8=5$\,Mbps per busy station.
  \item[(b)] Used $3\times5=15$\,Mbps; idle $40-15=25$\,Mbps (5 of every 8 slots are empty).
  \item[(c)] With $M=3$ busy stations, the ideal is $R/M=40/3\approx13.3$\,Mbps each.
\end{enumerate}
\ans{(a) $5$\,Mbps; (b) $25$\,Mbps wasted; (c) property 2 fails ($R/M$). Property 1 fails too: a station alone would still get only $5$\,Mbps, not $40$.}

\s{2.4}
\idea{Maximize by setting the derivative to zero (product rule), then let $N\to\infty$ using $(1-\frac1N)^N\to\frac1e$.}
(a) Slotted, $f(p)=Np(1-p)^{N-1}$:
\begin{enumerate}
  \item $f'(p)=N(1-p)^{N-1}+Np\cdot(N-1)(1-p)^{N-2}\cdot(-1)$.
  \item Factor out $N(1-p)^{N-2}$: $f'(p)=N(1-p)^{N-2}\big[(1-p)-(N-1)p\big]=N(1-p)^{N-2}\,[1-Np]$.
  \item $f'=0\Rightarrow p^*=1/N$ (the bracket is positive before and negative after, so it is a maximum).
  \item $f(p^*)=N\cdot\frac1N(1-\frac1N)^{N-1}=(1-\frac1N)^{N-1}=\dfrac{(1-\frac1N)^{N}}{1-\frac1N}\to\dfrac{1/e}{1}$.
\end{enumerate}
(b) Pure, $g(p)=Np(1-p)^{2(N-1)}$:
\begin{enumerate}
  \item In the same way: $g'(p)=N(1-p)^{2N-3}\big[(1-p)-2(N-1)p\big]=N(1-p)^{2N-3}\,[1-(2N-1)p]$, so $p^*=\frac1{2N-1}$.
  \item With $M=2N-1$: $g(p^*)=\dfrac{N}{2N-1}\Big(1-\dfrac1M\Big)^{M-1}$. The first factor $\to\frac12$ and the second $\to\frac1e$.
\end{enumerate}
\ans{(a) $p^*=1/N$, limit $1/e\approx0.37$; (b) $p^*=1/(2N-1)$, limit $1/(2e)\approx0.18$.}

\s{2.5}
\idea{First find the per-slot probabilities. Slots are independent, so ``first time in slot $k$'' means fail $k-1$ times, then succeed (geometric).}
With $p=0.25$: $(1-p)^3=0.4219$.
\begin{enumerate}
  \item[(a)] A wins one slot when A sends and B, C, D are silent: $q=p(1-p)^3=0.25\times0.4219=0.1055$.
  First time in slot 4: lose slots 1, 2, 3, win slot 4: $(1-q)^3q=0.8945^3\times0.1055=0.7157\times0.1055=0.0755$.
  \item[(b)] ``Some node succeeds'' $=$ exactly one of the four sends. The four events ``node $i$ alone'' are disjoint, so add them: $s=4p(1-p)^3=0.4219$.
  \item[(c)] Same idea with $s$: $(1-s)^3s=0.5781^3\times0.4219=0.1932\times0.4219=0.0815$.
  \item[(d)] Efficiency $=$ the probability that a slot is successful $=s$.
\end{enumerate}
\ans{(a) $(1-q)^3q=0.0755$; (b) $4p(1-p)^3=0.422$; (c) $(1-s)^3s=0.0815$; (d) $0.422$, the best possible for 4 nodes, since $p=1/4=1/N$.}
\watch{(a) is about one node ($q$, no factor 4); (b) and (c) are about any node ($s$, factor 4).}

\s{2.6}
\idea{A slot is A's if A sends and B does not. Write each node's success probability and compare.}
\begin{enumerate}
  \item[(a)] A: $p_A(1-p_B)$. B: $p_B(1-p_A)$. These are disjoint, so the efficiency is $p_A(1-p_B)+p_B(1-p_A)$.
  \item[(b)] Put $p_A=2p_B$: A gets $2p_B(1-p_B)$ and B gets $p_B(1-2p_B)$. Their ratio is $\dfrac{2(1-p_B)}{1-2p_B}$, which is more than 2 because $1-p_B>1-2p_B$.
  Reason: A sending more often also silences B more often.
  For exactly twice, solve $p_A(1-p_B)=2p_B(1-p_A)$: $p_A-p_Ap_B=2p_B-2p_Ap_B\Rightarrow p_A(1+p_B)=2p_B\Rightarrow p_A=\dfrac{2p_B}{1+p_B}$.
  Check with $p_B=0.1$: $p_A=0.182$; A gets $0.182\times0.9=0.164$, B gets $0.1\times0.818=0.082$, exactly half.
  \item[(c)] A wins: A sends ($2p$) and the other $N-1$ are silent: $2p(1-p)^{N-1}$. Another node $j$ wins: $j$ sends ($p$), A is silent ($1-2p$), and the remaining $N-2$ are silent: $p(1-2p)(1-p)^{N-2}$.
\end{enumerate}
\ans{(a) $p_A(1-p_B)$, total $p_A(1-p_B)+p_B(1-p_A)$; (b) no, more than twice; use $p_A=2p_B/(1+p_B)$; (c) $2p(1-p)^{N-1}$ and $p(1-2p)(1-p)^{N-2}$.}

\s{2.7}
\idea{Goodput $=$ P(exactly one sends); throughput $=$ P(at least one sends).}
\begin{enumerate}
  \item[(a)] $T=L/B=2000/(2\times10^6)=1$\,ms.
  \item[(b)] $\eta=\binom81\,0.1\cdot0.9^7=0.8\times0.4783=0.383$.
  \item[(c)] At least one sends $=1-P(\text{nobody})=1-0.9^8=1-0.4305=0.570$.
  \item[(d)] $\eta\times B=0.383\times2$\,Mbps $=765$\,kbps.
  \item[(e)] Successes are geometric with mean $1/\eta=2.61$ slots, i.e.\ $2.61\times1$\,ms.
\end{enumerate}
\ans{(a) $1$\,ms; (b) $0.383$; (c) $0.570$; (d) $765$\,kbps; (e) $\approx2.6$\,ms. The channel is busy 57\% of the time but useful only 38\%.}

\s{2.8}
\idea{Poisson model: the number of attempts in a slot is $X\sim\text{Pois}(G)$, so $P(X=k)=e^{-G}G^k/k!$.}
\begin{enumerate}
  \item[(a)] Success $P(X=1)=Ge^{-G}=0.4\times0.6703=0.268$. Empty $P(X=0)=e^{-0.4}=0.670$. Collision: the rest, $1-0.670-0.268=0.062$.
  \item[(b)] A pure frame succeeds if no other frame starts in its $2T$ window: $\text{Pois}(2G)$ is $0$ with probability $e^{-2G}$. So $\eta=Ge^{-2G}=0.4\,e^{-0.8}=0.4\times0.4493=0.180$.
  \item[(c)] Slotted is best at $G^*=1$ and pure at $G^*=0.5$. $G=0.4$ is below both.
\end{enumerate}
\ans{(a) $0.268$, $0.670$, $0.062$; (b) $0.180$; (c) increase the load in both cases: up to $G=1$ (slotted, max $0.368$) and slightly, to $G=0.5$ (pure, max $0.184$).}

\s{2.9}
\idea{(a) asks about \emph{C's} message: split on the network C chooses (total probability). (b) asks about the \emph{network}: a slot is useful if any one node succeeds.}
(a) One attempt.
\begin{enumerate}
  \item Upper (slotted): C sends ($q$) and each of the 3 others is silent in that one slot: $q(1-q)^3=0.1\times0.729=0.0729$.
  \item Lower (pure): C's frame overlaps \emph{two} slots of each other node, so each must skip both, $(1-q)^2$; with 2 nodes, $(1-q)^4$. So $q(1-q)^4=0.1\times0.6561=0.0656$.
  \item Combine: $p\times0.0729+(1-p)\times0.0656=0.6\times0.0729+0.4\times0.0656=0.0437+0.0262$.
\end{enumerate}
(b) Goodput of the upper network.
\begin{enumerate}
  \item C chose upper ($p$): 4 nodes, each sending with probability $q$; exactly one: $4q(1-q)^3=0.2916$.
  \item C chose lower ($1-p$): 3 nodes: $3q(1-q)^2=0.243$.
  \item $\eta=0.6\times0.2916+0.4\times0.243=0.175+0.097$.
\end{enumerate}
\ans{(a) $q[p(1-q)^3+(1-p)(1-q)^4]=0.070$; (b) $p\cdot4q(1-q)^3+(1-p)\cdot3q(1-q)^2=0.272$.}
\watch{one specific node (C) $\to$ no factor $n$; the whole network $\to$ factor $n$ ($\binom n1$).}

\s{2.10}
\idea{For each competitor, ask: how many of its slot starts fall in the window where a start would hit my frame? Each such start must be skipped.}
(a) A slotted node's slot is $[t,t+T]$.
\begin{enumerate}
  \item It sends: $p$.
  \item The other slotted nodes use the same aligned slot; each must be silent: $(1-p)^{n-1}$.
  \item A pure frame $[u,u+T]$ overlaps $[t,t+T]$ if $u\in(t-T,\,t+T)$, a window of length $2T$. A pure node starts every $T$, so exactly 2 of its starts fall there; it must skip both: $(1-p)^2$. With $m$ pure nodes: $(1-p)^{2m}$.
  \item Different nodes are independent, so multiply: $p(1-p)^{n-1}(1-p)^{2m}$.
\end{enumerate}
(b) A pure node's frame $[u,u+T]$, with $u$ between two grid points.
\begin{enumerate}
  \item It sends: $p$.
  \item The other pure nodes: 2 starts each in its window, as in plain pure ALOHA: $(1-p)^{2(m-1)}$.
  \item The frame covers the end of one grid slot and the start of the next, i.e.\ \textbf{two} slots. Every slotted node must be silent in both: $(1-p)^{2n}$.
  \item Multiply: $p(1-p)^{2(m-1)}(1-p)^{2n}$.
\end{enumerate}
(c) Slotted: $0.1\times0.9^{4}\times0.9^{10}=0.1\times0.9^{14}=0.0229$. Pure: $0.1\times0.9^{8}\times0.9^{10}=0.1\times0.9^{18}=0.0150$.
\ans{(a) $p(1-p)^{n-1+2m}$; (b) $p(1-p)^{2(m-1)+2n}$; (c) $0.0229$ vs.\ $0.0150$. The slotted node does better: it is exposed to one slot of each slotted node, while the unaligned pure frame is exposed to two slots of everyone.}

\s{2.11}
\idea{CSMA/CD only works if the sender is still transmitting when news of a collision comes back: $\TRANSP\ge2\,\PROP$.}
\begin{enumerate}
  \item $\PROP=d/s=200/(2\times10^8)=1\,\mu$s, so $2\,\PROP=2\,\mu$s.
  \item $\TRANSP=L/R\ge2\,\mu$s gives $L\ge2\times10^{-6}\times10^9=2000$ bits $=250$ bytes.
  \item Ethernet's 64 bytes $=512$ bits last only $512/10^9=0.512\,\mu$s $<2\,\mu$s.
\end{enumerate}
\ans{$250$ bytes. 64 bytes is too short at 1\,Gbps over 200\,m: a collision could go unnoticed. (In practice gigabit Ethernet is switched, so there are no collisions.)}

\s{2.12}
\begin{enumerate}
  \item After the $m$-th collision, $K$ is uniform on $\{0,\dots,2^m-1\}$. Here $m=5$: $\{0,\dots,31\}$, 32 values, so $P(K=4)=1/32$.
  \item The wait is $K\cdot512$ bit times $=2048$ bit times. One bit at 10\,Mbps takes $0.1\,\mu$s: $2048\times0.1=204.8\,\mu$s.
\end{enumerate}
\ans{$1/32$; $204.8\,\mu$s.}

\s{2.13}
\begin{enumerate}
  \item $K\cdot512=100\times512=51\,200$ bit times.
  \item 100\,Mbps: a bit time is $10$\,ns, so $51\,200\times10\,\text{ns}=512\,\mu$s. 1\,Gbps: $1$\,ns, so $51.2\,\mu$s.
\end{enumerate}
\ans{$512\,\mu$s and $51.2\,\mu$s.}

\s{2.14}
\idea{Draw the timeline in bit times. The worst case is B starting at the last possible moment, just before A's signal reaches it.}
\begin{enumerate}
  \item $t=0$: A starts. A's first bit reaches B at $t=325$.
  \item Just before $t=325$, B still senses an idle line and starts sending.
  \item B's first bit needs another 325 bit times to reach A: it arrives at $t\approx650$.
  \item A's minimum frame ends at $t=576$, before $650$.
\end{enumerate}
\ans{Yes: A finishes at 576, before hearing B at 650, so A wrongly believes its frame got through. The cable is too long for this frame size ($\TRANSP<2\,\PROP$).}

\s{2.15}
\idea{Two nodes collide again exactly when they pick the same $K$. With $k$ equally likely values, that happens with probability $k\cdot(1/k)^2=1/k$.}
\begin{enumerate}
  \item[(a)] After the 2nd collision: $\{0,1,2,3\}$.
  \item[(b)] 4 values: $1/4$.
  \item[(c)] To collide on the next three attempts they must match from sets of 4, then 8 (after the 3rd collision), then 16 (after the 4th). The choices are independent: $\frac14\cdot\frac18\cdot\frac1{16}$.
\end{enumerate}
\ans{(a) $\{0,\dots,3\}$; (b) $1/4$; (c) $1/512$. (This ignores the small effect of the busy-channel rule; it is the usual exam answer.)}

\s{2.16}
\idea{$\eta=\dfrac{\TRANSP}{\TRANSP+A\cdot2\,\PROP}$, where $A$ is the expected number of wasted $2\,\PROP$ slots per packet.}
\begin{enumerate}
  \item[(a)] $a=\PROP/\TRANSP=2/50=0.04$. Simple model ($A=1.5$): $\eta=50/(50+1.5\times4)=50/56=0.893$, which is the same as $1/(1+3a)$. Precise: $1/(1+5a)=1/1.2=0.833$.
  \item[(b)] $A=(1-\alpha)/\alpha$ with $\alpha=1/e$: $A=e-1=1.718$. Wasted time $1.718\times4=6.87\,\mu$s, so $\eta=50/56.87=0.879$.
\end{enumerate}
\ans{(a) $0.893$ (simple) and $0.833$ (precise); (b) $0.879$.}

\s{2.17}
\idea{The channel alternates between contention (wasted slots) and a frame ($k$ useful slots). Efficiency $=$ useful $\div$ total per cycle.}
\begin{enumerate}
  \item[(a)] A contention slot is won if exactly one node sends: $\alpha=Np(1-p)^{N-1}$. The number of wasted slots before a win is the number of failures before the first success of $\text{Geo}(\alpha)$: $x=(1-\alpha)/\alpha$. The winning slot is the first slot of the frame, so a cycle is $x$ wasted $+\,k$ useful slots: efficiency $=\dfrac{k}{k+(1-\alpha)/\alpha}$.
  \item[(b)] The efficiency grows as $x$ shrinks, i.e.\ as $\alpha$ grows. $\alpha$ is maximized at $p^*=1/N$ (Q2.4).
  \item[(c)] $\alpha\to1/e$, so $x\to(1-\frac1e)\,e=e-1\approx1.72$.
  \item[(d)] $\dfrac{k}{k+1.72}\to1$ as $k\to\infty$. For example $k=100$: $0.983$.
\end{enumerate}
\ans{(a) $\dfrac{k}{k+(1-\alpha)/\alpha}$ with $\alpha=Np(1-p)^{N-1}$; (b) $p^*=1/N$; (c) $\dfrac{k}{k+e-1}$; (d) $\to1$: long frames make contention negligible, the same as $a\to0$ in the lecture.}

\s{2.18}
\begin{enumerate}
  \item[(a)] Extra time per frame: $25+25=50\,\mu$s on a $500\,\mu$s frame, i.e.\ $10\%$.
  \item[(b)] Without RTS/CTS, a collision destroys a whole data frame: $500\,\mu$s. With it, only an RTS collides: $25\,\mu$s.
  \item[(c)] It costs $50\,\mu$s on every frame and saves time only when there is a collision.
\end{enumerate}
\ans{(a) $10\%$; (b) $500$ vs.\ $25\,\mu$s; (c) not worth it when collisions are rare (light load, no hidden terminals) or frames are short.}

\end{document}
